\scanpage{46}
schémas abéliens, \inserted{permettant de définir un foncteur qu.inverse du foncteur du théorème.} Pour ceci, on considère pour tout schéma abélien $A$ sur $S$ son extension vectorielle universelle, qui est une extension
\[
0\longrightarrow\underline t_{A^*}^{\vee}\;-\;E(A)\;-\;A\longrightarrow0.
\]
Il se trouve qu'on peut trouver un cristal en groupes lisses canonique sur $S$, soit $\underline E(\unsure{A})$, tel que l'on ait
\[
E(A)\unsure{\simeq|\cong}\underline E(A)_S.
\]
La définition est telle que \deleted{pour tout le foncteur}\inserted{$A$} $\underline E(A)$ définisse un foncteur
\[
\text{Schémab}(S)\longrightarrow\text{Cris Groupes lisses}_{\mathrm{nilp}}(S)
\]
compatible aux changements de base, l'isomorphisme précédent étant également fonctoriel et compatible aux changements de base. Il s'ensuit donc que l'on peut définir $\underline E(A)_{\deleted{S}}$ comme $S'\mapsto\underline E(A)_{S'}$, avec
\[
\underline E(A)_{S'}=E(A'),
\]
$A'$ désignant un schéma abélien qui prolonge $A$ sur $S'$, \deleted{\unsure{à}} à charge de définir \deleted{un}\inserted{un système transitif d'} isomorphismes canoniques entre les groupes obtenus pour des prolongements différents. Dans le cas où $S$ est de \unsure{caractééristique|caractéristique} nulle (donc plus question de puissances divisées !) cela est particulièrement simple; on prouve qu'il existe à isomorphisme \uline{unique} près un seul schéma en groupes lisses $E'$ sur l'épaississement $S'$ de $S$ qui prolonge $E$ (NB il suffirait même d'un nilépaississement, comme on voit par passage à la limite à partir du cas noethérien): c'est le résultat déjà signalé dans ma vieille lettre à Tate. Dans le cas général, c'est plus délicat, et peut se faire par ``la méthode de l'exponentielle'', qui sera exposée dans le contexte analogue des groupes de BT dans la suite du séminaire. On peut aussi définir $\underline E(A^*)$ par
\[
\underline E(A^*)=R^1f_{A\mathrm{cris}*}(\underline G_{m\,A_{\mathrm{cris}}}),
\]
où $A_{\mathrm{cris}}$ désigne le topos cristallin \deleted{\unsure{abso}} nilpotent absolu (relativement à $\underline Z$ ): cette méthode peut également s'adapter au cas des groupes de BT.

Quoi qu'il en soit de la méthode utilisée pour définir $\underline E(A)$, on
\margin{45}
